\section{An Independent Institutional Line}

\historythesis{The Institute of Christ the King Sovereign Priest did not arise by leaving the Society of Saint Pius X, and it was not formed as a later copy of the Priestly Fraternity of Saint Peter. Its recoverable line runs instead through Roman and Genoese formation, a bounded association with the French Opus Sacerdotale, diocesan recognition and missionary work at Mouila in Gabon, and a seminary and motherhouse at Gricigliano near Florence. The Institute made this geographically dispersed origin durable by joining priestly formation, solemn older Roman worship, common life modeled on canons, Salesian spirituality, diocesan invitations, and general government. Pontifical-right erection in 2008 and definitive constitutional approval reported in 2016 stabilized that form without removing either Roman supervision or the authority of bishops over local apostolates.}

The Institute is therefore not adequately described as a network of locations where the 1962 Missal is celebrated. It forms candidates, incorporates members, maintains a common rule and central government, assigns priests and non-priest oblates, and accepts pastoral works from bishops. The seminary and motherhouse at Gricigliano supply a common culture across countries. The associated Sister Adorers and the lay Society of the Sacred Heart share elements of spirituality while retaining identities that cannot be collapsed into the male society.

The reverse reduction is also unsound. The older Roman liturgy is not incidental to the Institute's self-understanding. The 2008 decree presents solemn liturgical worship according to the then-described extraordinary form as a defining service, and its account of common life places the Eucharistic sacrifice at the center of corporate and pastoral life. Later universal restrictions therefore affected more than scheduling, even though no checked act after 2021 suppressed the Institute or published a universal ICKSP-specific exemption.

\subsection{Question, range, and terminal date}

This account asks how Gilles Wach and Philippe Mora moved from Roman priestly formation into a new community; what can and cannot be established about the Opus Sacerdotale and the French passage; why Gabon and Florence became complementary institutional centers; how the Institute's ``canonial form'' works; and how its components, government, liturgical use, and apostolates changed from 1990 through 16 July 2026. It then examines the documentary record of pontifical-right erection, visitation, constitutions, local reception after \emph{Traditionis custodes}, papal contact, and publicly reported concerns.

TPI-01 supplies the longer history of the SSPX and its conflicts with Roman authority. TPI-02 supplies the FSSP's distinct branch point in 1988 and its special liturgical acts. Those volumes provide the wider field without making their institutions ancestors of the ICKSP. Cross-reference prevents a third retelling of Écône while preserving the differences among three bodies commonly grouped by worship.

The series title ``Traditional Priestly Institutes'' is bibliographic, not canonical. At the cutoff the ICKSP is a society of apostolic life of pontifical right. Its own phrase ``in canonial form'' describes a distinctive manner of common life; it does not by itself turn every house into a cathedral chapter or classify the society as an order of canons regular.

\subsection{Evidence layers and material absences}

\begin{threestable}{Layer}{Representative evidence}{Control and limit}
Public act & Papal records, Code of Canon Law, 2008 decrees, diocesan notices & Establishes the act and its printed terms. It does not supply facts absent from the text.\\
Institutional record & Histories, constitutions notice, directories, calendar, election notice, location pages & Establishes dated self-description. It is not an independent audit and sometimes contains internal chronological differences.\\
Historical reconstruction & Albert Jacquemin's history of the Opus Sacerdotale and bounded contextual synthesis & Helps distinguish an older priestly association from later institutional memory. It cannot replace missing erection acts.\\
Independent reporting & Reporting on Chicago and alleged governance or formation concerns & Establishes that attributed claims and responses were published, not that contested facts were adjudicated.\\
\end{threestable}

No original diocesan instruments for the 1990 Mouila erection or the early Florence establishment were located. The checked public set did not contain the 2014 visitation mandate or report, the full definitive constitutions or 2016 approval act, a post-2021 ICKSP-specific universal faculty, or a public Holy See outcome to the concerns reported in 2023 and 2026. Those absences are part of the history, not invitations to reconstruct private files.

\evidenceframe{Pontifical Commission \emph{Ecclesia Dei}, \emph{Saeculorum Rex}, N. 181/2008, 7 October 2008.}{The Holy See erected the Institute as a society of apostolic life of pontifical right, approved its constitutions experimentally for five years, described common life on the model of canons, and named Wach prior general for a renewable six-year term.}{The decree does not narrate every stage before 2008, enumerate the four 1962 books found in the FSSP's special acts, or erase the competencies of diocesan bishops.}
